\title{.25em{{.25em{ Paloma}}} : A Benchmark for Evaluating Language Model Fit}

\begin{abstract}
Evaluations of language models (LMs) commonly report perplexity on monolithic data held out from training. Implicitly or explicitly, this data is composed of domains—varying distributions of language. We introduce \textsc{Perplexity Analysis for Language Model Assessment} (\textsc{Paloma})\footnote{Dataset and links to code repository are available at \url{https://paloma.allen.ai}}, a benchmark to measure LM fit to 546 English and code domains, instead of assuming perplexity on one distribution extrapolates to others. We include two new datasets of the top 100 subreddits (e.g., \textit{r/depression} on Reddit) and programming languages (e.g., \textit{Java} on GitHub), both sources common in contemporary LMs. With our benchmark, we release 6 baseline 1B LMs carefully controlled to provide fair comparisons about which pretraining corpus is best and code for others to apply those controls to their own experiments. Our case studies demonstrate how the fine-grained results from \textsc{Paloma} surface findings such as that models pretrained without data beyond Common Crawl exhibit anomalous gaps in LM fit to many domains or that loss is dominated by the most frequently occurring strings in the vocabulary.
\end{abstract}

\section{Introduction}

Progress in AI is catalyzed by evaluations that define new ways of measuring progress (\citealp{imagenet}, \citealp{wang-etal-2018-glue}, and \citealp{superglue}, \textit{inter alia}). Language models (LMs) often evaluate LM fit as loss or perplexity \citep{Jelinek1977PerplexityaMO} on held out training data or few traditional test sets (\citealp{Chelba2013OneBW,Merity2016PointerSM}, \textit{inter alia}). These loss measures have been shown to improve predictably with increases in training compute \citep{Kaplan2020ScalingLF, Hoffmann2022TrainingCL} and loss may predict performance on downstream tasks \citep{Xia2022TrainingTO, Gadre2024LanguageMS, Du2024UnderstandingEA}. However, scaling pretraining data aggregates more domains that LMs implicitly learn to model \citep{Diaz2023ScalingLD, aharoni-goldberg-2020-unsupervised}. Does rising performance lift all data? Or do some domains capture most improvement in LM fit? How do we evaluate what language distributions models learn from different pretraining data? What domains should studies evaluate loss on to measure the relationship of loss and downstream performance? To answer these questions, perplexity evaluations ought to measure LM fit to many domains, rather than extrapolating trends from a single prescriptive mix of domains. 

In this work we introduce \textsc{Paloma}, a benchmark to study LM fit on many domains. We measure perplexity on different distributions of language sampled from 16 sources, such as \textsc{C4} \citep{Raffel2019ExploringTL}, that have metadata such as URLs marking 546 textual domains. Beyond evaluation data, we aim to enable and enrich fair comparisons for scientific research on language modeling with the following artifacts: guidelines for comparing LM fit, 6 baseline 1B parameter models pretrained on popular corpora, and standardized code for experiments with \textsc{Paloma}.

As reproducing pretrained models for every new project is onerous, we provide standard training controls for benchmark decontamination and training data order to orchestrate a greater density of comparisons across the research community. We also control how \textsc{Paloma} is evaluated by fixing sample size per domain, model vocabulary, and inference format. Lastly, we demonstrate how to make fair comparisons over two measures of cost, number of model parameters and training tokens, enabling assessment of hardware-agnostic efficiency and the measurement of scaling trends.

Among the 16 sources curated in our benchmark, we contribute two new datasets constructed from data held out of \textsc{Dolma} \citep{dolma}: (1) a subsample of the top 100 subreddits by number of comments, and (2) code from the top 100 programming languages by number of tokens. Also, we repurpose corpora of fringe online communities to measure LM fit to discourse previously studied for the prevalence of toxicity and hate speech \citep{Horta_Ribeiro_2021,ZannettouGab,Papasavva_2020}. While, capturing domains required by all possible lines of research is impossible for any one benchmark, \textsc{Paloma} focuses on English and code data and aims to assemble the most fine-grained domains readily identifiable from existing metadata.

To demonstrate possible uses of results from our dataset, we present a series of case studies in \S\ref{sec:case_studies}. 

Among other findings, our experiments isolate change in fit from which pretraining corpus is used (Figure~\ref{fig:models_by_ppl_over_all_tasks}) and find that pretraining without heterogeneous data sources beyond Common Crawl can lead to perplexities in some domains that do not improve consistently with number of tokens seen. We also find that few vocabulary types account for most of the loss measured in perplexity.

In sum, \textsc{Paloma} contributes:

\begin{enumerate}
    \item Curated release of the most fine-grained perplexity evaluation data in use in LM research, along with guidelines and code for standardized and rigorous perplexity evaluation.
    \item New evaluation data for the 100 most popular subreddits and programming languages.
    \item 1B LMs pretrained on \textsc{C4}, \textsc{mC4-en}, \textsc{Falcon Refinedweb}, \textsc{The Pile}, \textsc{RedPajama}, and \textsc{Dolma} with controlled hyperparameters, token budget, benchmark decontamination, and training order for fair comparisons, along with code for others to do the same.
    \item Case studies demonstrating analyses that are possible with \textsc{Paloma}, such as finding that pretraining without data beyond Common Crawl leads to inconsistent fit to many domains and that perplexity is driven by improved fit on the most common vocabulary strings.
\end{enumerate}

\section{Sources of evaluation data}
\label{sec:selecting}

We define two terms: \emph{Sources} are as existing datasets (or curated subsets there of) in use for research. \emph{Domains} are fine-grained partitions of sources based on available metadata that attempt to surface a distinct and intuitive distribution of language (e.g., Wikipedia articles about visual arts or a subreddit for advice on PC builds). \textsc{Paloma} is derived from 16 sources further divided into 546 domains (see Table~\ref{tab:source_statistics}).\footnote{Unless stated, token counts are computed with the GPT-NeoX-20B tokenizer \citep{gpt-neo}.} Where we curate previous fine-grained corpora, we inherit their operationalization of domains, ranging from the community-driven Wikipedia ontology to expert curation and automatic classification. Where we build our own fine-grained domains from Reddit and GitHub, we make similar use of metadata about subreddits and file extensions.

Compared to monitoring monolithic validation loss during model development, interpreting LM fit to specific fine-grained domains poses unique challenges. Crucially, we must not assume better LM fit to a domain reflects improvements in the specific skills that are valued by the humans producing language in that domain \citep{Diaz2023ScalingLD}. For instance, we might expect overlapping domains for academic papers in both \textsc{Dolma} and \textsc{RedPajama} to exhibit similar perplexities for a given model, perhaps assuming perplexity represents how much a model captures knowledge about relevant academic fields. But domains can also differ due to preprocessing when texts were collected in each source rather than from how texts were composed by their original authors. So instead of relying on LM fit to measure what we think a model \textit{should} learn about a domain, we examine anomalies in domain fit to see what a model \textit{is} learning. We find that the same model can have 391,171 perplexity on arXiv in \textsc{RedPajama} and 14 on the overlapping academic domain, peS2o, in \textsc{Dolma} (\S\ref{sec:isolating_pretraining}). In this approach we follow \citet{Holtzman2023GenerativeMA} and \citet{McCoy2023EmbersOA} by aiming to examine model behaviors, regardless of their desirability to humans.

Also note that \textsc{Paloma} focuses on English and code data, as most current LMs also emphasize these types of data. However, we strongly encourage future work to explore fit to fine-grained domains in other languages. 

The rest of this section addresses each source, why we include it, and how it identifies any domains it contains (all 546 domains are listed in Appendix~\ref{sec:source_details}).

\paragraph{Standard language modeling sources} Though it is common practice to evaluate on held out data from the pretraining corpus of a given model, we evaluate \textit{across} several standard corpora. \textbf{\textsc{C4}} \citep{Raffel2019ExploringTL, dodge-etal-2021-documenting} and \textbf{\textsc{mC4-en}} \citep{Chung2023UniMaxFA} are language model training datasets created by taking the snapshots of Common Crawl data and applying a number of filters with the intention of retaining ``high-quality'', natural language. Both datasets are filtered to retain natural English, and in this work we only use the English portion of \textsc{mC4-en}.
\textbf{\textsc{WikiText-103}} \citep{Merity2016PointerSM} and \textbf{\textsc{Penn Treebank}} \citep{MarcusPTB} are classic datasets that have been used to evaluate language model perplexity for decades (\citealp{Radford2019LanguageMA, brown2020language, Rae2021ScalingLM,Hoffmann2022TrainingCL}, \textit{inter alia)}. \textsc{WikiText-103} is text from Wikipedia articles, and \textsc{Penn Treebank} \citep{MarcusPTB} is a set of 1989 Wall Street Journal articles\footnote{\textsc{Penn Treebank} is pretokenized, and uncommon words are replaced with a special ``unknown'' token.}.
\textbf{\textsc{RedPajama}} \citep{together2023redpajama} is an attempt at reproducing the data mixture from LLaMA \citep{llama} from sources such as webtext, Wikipedia, arXiv, and StackExchange. It was used to train RedPajama-INCITE \citep{together2023redpajama}.
\textbf{\textsc{Falcon Refinedweb}} \citet{Penedo2023TheRD} was created from all Common Crawl scrapes until June 2023  by applying relatively interpretable filters, and is a subset of the Falcon models' training data \citep{falcon}.
\textbf{\textsc{Dolma}} \citet{dolma} is made of Common Crawl, Wikipedia, books, academic papers, code repositories, and Reddit, and was used to train OLMo models \citep{Groeneveld2024OLMoAT}.

\paragraph{Fine-grained domain sources} We include datasets with the most fine-grained metadata marking hundreds of domains.
\textbf{\textsc{M2D2}} \citep{reid-etal-2022-m2d2} is made of academic papers from S2ORC \citep{lo-etal-2020-s2orc} and text from Wikipedia, organized into a two-level hierarchy by academic field categories or Wikipedia ontology, respectively. We sample both top-level domains and lower-level subdomains. 
\textbf{\textsc{C4-100-domains}} \citep{chronopoulou-etal-2022-efficient} is text from the 100 internet domains with the most pages in \textsc{C4}.\footnote{Four of the 100 domains have less than the 100 thousand tokens per split that we aim for.}
\textbf{\textsc{Dolma-100-subreddits}} and \textbf{\textsc{Dolma-100-programming-languages}} are two evaluation sets we introduce in this work sampled from \textsc{Dolma} \citep{dolma}: the former is text from the top 100 subreddits (ranked by number of posts), and the latter is the top 100 programming languages by number of tokens in the \textsc{The Stack} \citep{Kocetkov2022TheStack}. See Appendix~\ref{sec:source_details} for more details.

\paragraph{Disparities between speech communities} LMs today primarily process dominant dialects in countries, such as the US, where they are most often trained and deployed. Even within English, hundreds of millions of people around the world speak other dialects that have been shown to be underserved by existing models \citep{blodgett-etal-2016-demographic}. As a starting point for measuring disparities between dialects, we include
\textbf{\textsc{TwitterAAE}} \citep{blodgett-etal-2016-demographic}, two corpora representing African-American and White-aligned English, automatically classified via geolocation information and demographic census statistics. \footnote{We follow the reproduction of this dataset used in HELM \citep{Liang2022HolisticEO}, but we fix an error in loading escaped sequences of the data that, among other issues, renders emojis as literal hexadecimal bytes.}

\paragraph{Fringe sources previously studied for problematic discourse} LM fit to these fringe texts characterizes model exposure to distinct social contexts in which toxic language arises.
\textbf{\textsc{Manosphere}} \citep{Horta_Ribeiro_2021}, \textbf{\textsc{Gab}} \citep{ZannettouGab}, and \textbf{\textsc{4chan Corpora}} \citep{Papasavva_2020} are three fringe corpora which contain larger proportions of hate speech and toxicity than mainstream sources like Wikipedia or Twitter. These texts span 2006-2019 and include independent message boards and subreddits sharing a masculinist ideology, Gab (an alt-right focused Twitter alternative with minimal moderation), and the Politically Incorrect board (/pol/) of 4chan, a fringe imageboard emphasizing anonymity and ephemerality.

\section{Perplexity evaluations done right}
\label{sec:guidelines}
\label{sec:submission_tracks} 

\paragraph{Guidelines}
Fairly evaluating different models using perplexity is hard. To do so, we must account for factors that can confound results with guidelines for training (\textsc{G1}, \textsc{G2}) and evaluation (\textsc{G3}, \textsc{G4}, \textsc{G5}).

\begin{enumerate}[label=\textsc{G\arabic*}]
    \item \textsc{Decontamination}: Remove \textit{pretraining} data that leaks evaluation data to ensure validity of perplexity evaluation. 
    \item \textsc{Training Order}: Where possible, keep the training data order the same to control differences from recency effects.
    \item \textsc{Subsampling}: Subsample size poses a tradeoff between inference cost and variance. Size subsamples to tolerate variance equally for each domain.
    \item \textsc{Vocabulary}: Vocabulary determines the event space of possible sequences and the comparability of perplexity measurements. Normalizing likelihood by a segmentation intrinsic to the text (e.g., bytes) partially addresses this, but fixing the vocabulary is preferable.
    \item  \textsc{Evaluation Format}: Use a consistent implementation of perplexity to ensure comparability regarding engineering details such as the handling maximum sequence lengths.
\end{enumerate}

\label{sec:experimental_controls}
\paragraph{Experimental controls} Our code repository\footnote{\url{https://github.com/allenai/OLMo-Eval/tree/main/paloma}} releases controls that implement each guideline. Here we briefly explain each (complete specification of our experimental controls is provided in Appendix~\ref{sec:controls}).

For \textsc{G1}, we use a Bloom filter \citep{Bloom1970SpacetimeTI} to detect exact match overlaps of pretraining and evaluation data. We match text at the paragraph level, i.e., newline separated spans of text. To avoid coincidental collisions in the space of small strings, we ignore matches in paragraphs smaller than 13 unicode segmented tokens \citep{uniseg}. Similarly, we ignore paragraphs composed of only punctuation, spaces, and emoji. Lastly, as code data consists almost entirely of short and often repeated lines, we forgo any decontamination on these sources (\textsc{Dolma-100-programming-languages} and the \textsc{The Stack} domain of \textsc{Dolma}). Finally, we remove whole pretraining documents if they contain \textit{any} contaminated paragraph.

For \textsc{G2}, contemporary LMs train on instances that are maximum sequence length concatenations of training documents, so we must fix the order of concatenated instances. We achieve this by fixing the tokenization, maximum sequence length, and random seed, as well as providing dataloading code where order is invariant to number of devices.

For \textsc{G3}, we empirically observe how variance in perplexity over subsamples of \textsc{C4} evaluation data grows inversely to sample size (Appendix~\ref{sec:subsampling_control}). Extrapolating from these results to select desired thresholds for variance, we pick 1 million and 100 thousand tokens as our target size for sources and domains, respectively.

For \textsc{G4}, where possible we fix model vocabulary to GPT-NeoX-20B's \citep{gpt-neo} with 3 special tokens added by \citet{Groeneveld2024OLMoAT}. When vocabulary must be changed, for instance comparing to off-the-shelf models, we follow \textsc{The Pile} \citep{Gao2020ThePA} and use bits per byte (BPB; Appendix~\ref{sec:additional_metrics}).

For \textsc{G5}, we follow the input format established by \textsc{The Pile} \citep{Gao2020ThePA}. This format evaluates documents individually, rather than packed into concatenated maximum sequence length inputs. Documents longer than maximum sequence length are split into disjoint inputs. 

In Table~\ref{tab:benchmark_diffs} we compare how \textsc{Paloma} implements controls for these guidelines against practices in previous LM benchmarks. \textsc{Paloma} is the first benchmark to remove contamination across all pretraining data. \textsc{The Pile}~\citep{Gao2020ThePA} note that they only address decontamination partially by deduplicating 2 of 22 domains at the document level before splitting. \textsc{Paloma} is also the first contemporary perplexity benchmark to recommend and implement a method to fix the training data order, to apply stratified sampling to evaluation domains, and to recommend fixing vocabulary. \textsc{The Pile} and HELM also detail their evaluation formats, but we note that HELM's inference code depends on calls to proprietary APIs which may not remain reproducible for some models.

\paragraph{Comparability} When using \textsc{Paloma} to compare models, we recommend that researchers also adopt our experimental controls or note as a limitation to comparability any uncontrolled factors. We also recommend that measures of cost are considered when comparing models on \textsc{Paloma}, specifically number of model parameters and number of tokens seen in training. Complimentary to work that focuses on realized costs such as energy use, FLOPs, or GPU hours \citep{Peng2023EfficiencyPA}, we elect to measure these more abstract cost values so that our efficiency comparisons are agnostic to hardware. Finally, as LMs trained with non-constant learning rate schedules scale sub-optimally until improving when learning rate drops towards the end of training, fair comparisons involving intermediate checkpoints should be matched with respect to the portion of total optimization steps completed.

By providing fair comparisons, the following types of claims about perplexity performance can be made with our benchmark: (1) which among compute-matched models performs best, (2) which models reach a given performance with the least compute, (3) which pretraining corpus produces models with best performance, (4) quantifying the trend of performance as a function of scale.

\label{sec:metrics}
\paragraph{Metric} \textsc{Paloma} uses standardized inference code to compute metrics to assess LM fit to the evaluation data we have curated. Perplexity \citep{Jelinek1977PerplexityaMO} is our primary metric (others not used in the body of this paper are detailed in Appendix~\ref{sec:additional_metrics}). Unless otherwise stated, we use perplexity to mean perplexity per token, where a log likelihood $\ell$ over documents $N = \{ t^{1},\ldots,t^{| N |} \}$ is normalized by $\mathbf{T}({N})$ denoting the number of tokens in the documents (i.e., $\mathbf{T}(N) = \sum_{t \in N} | \mathbf{tokenize}(t) |$):

\begin{myequation}
    \ell = \sum _{t \in N} \sum _{i}^{ | t |} \text{ln}~p(t_i \mid t_{<i})
\end{myequation}

\begin{myequation}
    \text{perplexity} = e^{-\frac{\ell}{\mathbf{T}(N)}}
\end{myequation}

\section{Case studies}
\label{sec:case_studies}

In this section, we present one full case study and a single conclusion from a second. In Appendix~\ref{sec:additional_case_studies} we present additional studies, demonstrating the types of analyses possible with \textsc{Paloma}.
\subsection{Pretraining Beyond Common Crawl Shows Improved Stability of LM Fit}
\label{sec:isolating_pretraining}
We hypothesize that one of the strongest drivers of differences in performance between different domains is the composition of the pretraining data of a language model. While we show in Appendix~\ref{sec:fine_grained_domain_tasks} that scaling model parameters or tokens seen increases performance on nearly all domains, the pretraining data composition directly determines the distribution of language that the model is learning to fit, which may or may not align with the distributions of language in the domains we evaluate. Therefore we examine the impact of varying the pretraining corpus while holding all other experimental decisions the same.

\label{sec:baseline_models}
\paragraph{Baseline Models} We train and release a set of 6 baseline models on common pretraining corpora following our training guidelines (\S\ref{sec:guidelines}). Training these models ourselves allows us to apply decontamination and fixed order to their pretraining data as well as using a standard tokenizer to enable the greatest level of comparability. These models are 1B parameter models trained for $\sim$150B tokens on \textsc{Dolma} \citep{dolma}, \textsc{The Pile} \citep{Gao2020ThePA}, \textsc{RedPajama} \citep{together2023redpajama}, \textsc{Falcon Refinedweb} \citep{Penedo2023TheRD}, \textsc{C4} \citep{Raffel2019ExploringTL,dodge-etal-2021-documenting}, and \textsc{mC4-en} \citep{Chung2023UniMaxFA}. Additional training details are included in Appendix~\ref{sec:baseline_models_appendix}.

\paragraph{Ordinary perplexity} In Figure~\ref{fig:models_by_ppl_over_all_tasks}, we consider the most simple and aggregated view of LM fit that \textsc{Paloma} can provide—perplexity as defined in \S\ref{sec:metrics}. Specifically we compute perplexity over all data, excluding the three fringe sources with prevalent toxicity. We also exclude code data in \textsc{Dolma} and \textsc{Dolma-100-programming-languages}.\footnote{We do not decontaminate code as its paragraphs (lines) are short and often repeated.}

Using this view, we see that baseline models trained only on Common Crawl data (\textsc{C4}, \textsc{Falcon Refinedweb}, and \textsc{mC4-en}) stand out from the others which incorporate more curated data sources. However, this points to the limitation of this most aggregated view of the results: ordinary perplexity represents fit to domains in proportion to the number of tokens we have chosen to sample from each domain. We sample 100,000 tokens from each domain and the majority of our domains are not sourced from Common Crawl. So Common Crawl is much less represented in \textsc{Paloma} than in most pretraining corpora, which typically consist of mostly Common Crawl as this is the most abundant public source of text data. Nevertheless this simplified view of the results is useful for specific use cases that need a single metric over a prescriptive mix that emphasizes robustness to a diversity of domains, largely derived from non-web scraped sources.

\paragraph{Macro average perplexity} Figure~\ref{fig:models_by_macro_subdomains}
provides another aggregation that examines the robustness of fit by considering all domains equally—a macro average of perplexity over domains: $|D|^{-1}\sum_{d\in D} \text{perplexity}(d)$ for domain set $D$. By contrast \textit{ordinary perplexity} is essentially an exponentiated micro average over the domains implicitly selected for during corpus curation. Macro averaging lets all marked domains have equal say on the model's performance, instead. To make these macro averages more easily interpretable, we examine them separately per source.

The most striking pattern that emerges with per-source macro averages is the high, and sometimes non-monotonic, perplexity of the 3 baselines trained on only Common Crawl data (\textsc{C4}, \textsc{mC4-en}, \textsc{Falcon Refinedweb}). This is particularly apparent for the \textsc{C4} model evaluated on \textsc{RedPajama}, where the macro average is dominated by perplexity up to 391,171 on the \textit{arXiv} domain. Similar spikes occur for the \textsc{Falcon Refinedweb} and \textsc{mC4-en} models, with perplexity of 21,652 and 1,409 respectively, on the Max music programming language domain in \textsc{Dolma-100-programming-languages}. These domains contain large amounts of non-natural language, in the form of LaTeX and other code data. These spikes stand out from the stable and monotonic improvement observed in the other 3 baseline models. While these Common Crawl baselines spike on different domains, it appears they are more susceptible to these extreme gaps in fit to \textit{some} domains. Perhaps this occurs because of a lack of exposure to specific types of language completely filtered due to having only one set of cleaning filters applied to a single source of data.

In contrast, the baselines that include curated non-webscraped text sources (\textsc{Dolma}, \textsc{The Pile}, and \textsc{RedPajama}) have a relative gap in perplexity that is highly stable through the course of training. This would imply that short training runs on a subsample of such pretraining corpora may be predictive of the LM fit of specific sources after much longer training. To address one exception, the \textsc{RedPajama} baseline often spikes on its final checkpoint, sometimes dramatically as in \textsc{TwitterAAE}. A possible explanation is that this checkpoint falls very soon after the model's training loss recovers from a small spike.

\paragraph{Perplexity per domain ordered by median perplexity} We can visualize each perplexity separately for each domain to surface gaps in fine-grained LM fit. In Figure~\ref{fig:subdomain_perf_ordered_by_task}, we arrange the domains by their median perplexity over the baselines, as this order gives some sense of the intrinsic difficulty of a domain. We can then see which baselines follow this order, differing only by a consistent offset, and which have gaps that are more idiosyncratic to each domain. Again we see that when baselines have irregular gaps from the median these are most frequently baselines pretrained on only Common Crawl. The notable exception is \textsc{The Pile} baseline on \textsc{M2D2 S2ORC} and \textsc{Dolma-100-programming-languages}, which has erratic gaps substantially below the median, perhaps indicating that baseline is benefiting from exposure to specific domains and not others rather than only a overall facility for scientific papers and code. The erratic-gapped Common Crawl baselines, by contrast, are all worse than median perplexity, suggesting that they may have complete gaps in exposure to features of certain domains that are not recovered through generalization.

\subsection{Common Vocabulary Types Dominate Perplexity}
\label{sec:common_types_body}
Here we present a single conclusion from a second case study; see Appendix~\ref{sec:token_level} for further analysis. So far we have examined perplexity aggregated over tokens. Another approach is to measure average likelihood per vocabulary \textit{type}, i.e., the strings that are represented in the vocabulary of a model, in contrast to occurrences of these strings in some corpus, called \textit{tokens}.\footnote{See Appendix~\ref{sec:additional_metrics} for a more formal definition of average likelihood per vocabulary type}

\paragraph{Few vocabulary types account for most of the loss measured in perplexity} How much do specific \emph{types} contribute to perplexity aggregated per token? To answer, we start by analyzing the total loss mass added by types, as a function of their IDs. Smaller IDs correspond to more frequent types in the GPTNeoX-20B tokenizer training data \citep{sennrich-etal-2016-neural, gpt-neo}, and we find an overall moderate to strong correlation between IDs and frequencies in the evaluation data of \textsc{Paloma} as well (Pearson's $r$ averaged across domains: \mbox{--0.522}$\pm$\mbox{0.087}). Crucially, frequency has a strong impact on the total loss mass associated with individual types: while the \emph{average} loss is lower for the high-frequency types (Figure~\ref{fig:mean-loss}), the \emph{total} loss is higher, resulting in a situation where 5\% of the types already cover roughly 50\% of the overall perplexity (Figure~\ref{fig:total-loss-cdf}). Thus, perplexity is strongly influenced by a relatively small set of high-frequency types. This finding provides further evidence that reporting only aggregated perplexity values neglects more subtle dynamics visible through fine-grained analysis (i.e., sources, domains, vocabulary types) in \textsc{Paloma}.

\section{Conclusion}
We believe that evaluations of language modeling fit provide an important view of performance that has been neglected in recent LM research and development.

Perplexity cannot be na\"{i}vely applied to language modeling at this scale due to challenges such as benchmark contamination. However, these obstacles are worth overcoming as perplexity offers several advantages not afforded by downstream evaluations. Instead of constructing tasks from scratch, we can rely on the ecological validity of real-world data drawn from known sources. Finding the best ways to evaluate model fit to a collection of documents creates an interface for other fields to contribute to the evaluation of language models. Without needing to understand LM architectures, researchers in other fields can collect corpora representing domains of interest that LM researchers would not know to consider. Once such sources are identified, evaluations can be updated over time by simply scraping more data, unlike downstream tasks where expensive annotation would be required. 

Further, we hope that \textsc{Paloma} provides controlled results for study of when perplexity evaluations are or are not predictive of downstream performance \citep{Liu2022SamePL, Tay2021ScaleEI, Ganguli2022PredictabilityAS,Xia2022TrainingTO, Gadre2024LanguageMS, Du2024UnderstandingEA}. In Appendix~\ref{sec:downstream_analysis}, our preliminary investigation reveals that different \textsc{Paloma} sources are correlated with some downstream tasks and anticorrelated with others. This contrasts with the assumption in much scaling literature that lower perplexity always indicates better downstream performance. While we do observe that LM loss reduces with scale across most domains (Appendix~\ref{sec:fine_grained_domain_tasks}), the fit of this relationship and the relationship of loss to downstream performance will both differ for each pretraining and validation distribution as observed by \citet{Gadre2024LanguageMS}. This means that one cannot simply find which fine-grained perplexity domains correlate with one's favorite task and then hillclimb on those. Instead further investigation with pretraining experiments across a wide range of scales and data recipes is needed to understand when reductions in perplexity are being driven by superficial overlaps of train and validation distributions or by learning features relevant to downstream use.

\section{Limitations and Future Work}
\label{sec:limitations}

The largest limitation of \textsc{Paloma} is that we elect to focus just on the language modeling of English and code data. We select this scope as most current LMs also focus on theses types of data. However, we strongly encourage future work to explore how language model fit to fine-grained domains behaves within and across other languages. 

Proper use of perplexity as a metric must take into account its limitations. We believe perplexity is best used to show what a model \textit{is} learning rather than what it \textit{should} be learning. For instance we find that perplexity on the 3 fringe datasets are tightly related to average document lengths, with the short tweet-like posts in \textsc{Gab Corpus} receiving high perplexities while the long concatenated threads of posts in \textsc{4chan Corpus} and \textsc{Manosphere Corpus} provide greater context and lower perplexity. At this level of aggregation, differences in surprise between these domains likely have little to do with model fit to specific types of toxicity and more to do with how models use extremely short or long contexts. In our case study in \S\ref{sec:common_types_body}, we demonstrate that often it is more appropriate to decompose measures of surprise over specific strings within a corpus, rather than aggregating over all text in a domain. We hope that by surfacing the average likelihoods of specific strings in the vocabulary, \textsc{Paloma} can enable future work on metrics that better measure the fit of models to the features of language in specific domains that humans find most salient.

We also highlight guidelines for evaluating with perplexity (\S\ref{sec:guidelines}). In particular we believe decontamination of benchmark leakage and balancing variance induced by subsampling across domains are both challenging concerns requiring further investigation. For each of these we have proposed one simple and scalable mitigation (see Appendix~\ref{sec:decontamination_control} and \ref{sec:subsampling_control} for further details), but future work should explore alternatives and measure their efficacy.

\textsc{Paloma} curates and standardizes the text datasets with the most fine-grained domains readily available from existing metadata. As such, our definition of domains by metadata is necessarily heuristic. Some overlapping domains in \textsc{Paloma} appear in multiple sources, such as academic papers. Though \textsc{Dolma} and \textsc{RedPajama} process academic papers differently, the subcorpora on academic papers in each source represent different approximations of the same or very similar domains. However for the sake of simplicity, we make the reductive assumption of counting all 546 domains in \textsc{Paloma} as fully distinct. We hope that future work will explore novel means of identifying fine-grained domains and separating distribution shifts in language due to differing authorship or differing data collection processes.

\end{document}
